\documentclass[12pt,a4paper]{amsart}

\usepackage{amsmath,amssymb,amsthm}
\usepackage{mathtools}
\usepackage[utf8]{inputenc}
\usepackage[T1]{fontenc}
\usepackage{hyperref}
\usepackage{enumitem,booktabs,array}
\usepackage{geometry}
\geometry{margin=2.5cm}

% -------------------------------------------------------
% Theorem environments
% -------------------------------------------------------
\newtheorem{theorem}{Theorem}[section]
\newtheorem{lemma}[theorem]{Lemma}
\newtheorem{corollary}[theorem]{Corollary}
\newtheorem{proposition}[theorem]{Proposition}
\theoremstyle{definition}
\newtheorem{definition}[theorem]{Definition}
\newtheorem{example}[theorem]{Example}
\newtheorem{remark}[theorem]{Remark}
\newtheorem{conjecture}[theorem]{Conjecture}

% -------------------------------------------------------
% Custom commands
% -------------------------------------------------------
\newcommand{\N}{\mathbb{N}}
\newcommand{\Z}{\mathbb{Z}}
\newcommand{\Q}{\mathbb{Q}}
\newcommand{\R}{\mathbb{R}}
\newcommand{\C}{\mathbb{C}}
\DeclareMathOperator{\ord}{ord}
\DeclareMathOperator{\lcm}{lcm}
\DeclareMathOperator{\vp}{v_p}

% -------------------------------------------------------
\begin{document}
% -------------------------------------------------------

\title{On the Four-Prime Case of Giuga's Conjecture:\\
       The Equation System, Structural Constraints,\\
       and Computational Evidence}

\author{Michael Fuhrmann}
\address{Specialist Mathematics Project, Build 133}
\email{specialist-maths@localhost}
\date{\today}

\begin{abstract}
Giuga's 1950 primality conjecture asserts that no composite integer satisfies
$\sum_{k=1}^{n-1} k^{n-1} \equiv -1 \pmod{n}$; equivalently, that no
\emph{Giuga pseudoprime} exists.  It is known that every such pseudoprime must
be squarefree and have at least four distinct prime factors, the three-prime
case having been ruled out by elementary methods.

This paper investigates the four-prime case in depth.  Writing
$n = p_1 p_2 p_3 p_4$ with $p_1 < p_2 < p_3 < p_4$ prime, we derive the
complete system of divisibility conditions that any such $n$ must satisfy
and analyse the resulting arithmetic constraints.  We prove the totient
identity $\varphi(n) = n - e_3 + e_2 - e_1 + 1 > 0$, showing that the
expression $e_4 - e_3 + e_2 - e_1 + 1 = 0$ is algebraically impossible
for any product of four distinct primes.  We establish several necessary
arithmetic constraints, including the sharp bound $p_4 \le p_1 p_2 p_3 - 1$
from condition~(G4) and the size estimate $n < (p_1 p_2 p_3)^2$.  We further
report on a computational search covering all products $n = p_1 p_2 p_3 p_4
\le 10^{15}$, finding no Giuga pseudoprime in this range.  The four-prime
case of Giuga's conjecture remains open.
\end{abstract}

\maketitle
\tableofcontents

% -------------------------------------------------------
\section{Introduction}
\label{sec:intro}
% -------------------------------------------------------

\subsection{Historical Background}

In 1950, the Italian mathematician Giuseppe Giuga posed a deceptively simple
question about a characterisation of prime numbers~\cite{Giuga1950}.
He observed that for every prime $p$ and every $k$ with $1 \le k \le p-1$,
Fermat's Little Theorem gives $k^{p-1} \equiv 1 \pmod{p}$, so
\[
  \sum_{k=1}^{p-1} k^{p-1} \equiv p - 1 \equiv -1 \pmod{p}.
\]
Giuga asked whether the converse holds: is every integer $n > 1$ satisfying
$\sum_{k=1}^{n-1} k^{n-1} \equiv -1 \pmod{n}$ necessarily prime?

A composite number satisfying this congruence would be called a
\emph{Giuga pseudoprime}.  Despite intensive study over more than seven
decades, no Giuga pseudoprime has ever been found, and the existence question
remains one of the outstanding open problems in elementary number theory.

\subsection{The Algebraic Characterisation}

A fundamental result, established by Borwein, Borwein, Girgensohn, and
Parnes~\cite{Borwein1996}, reduces Giuga's conjecture to a purely divisibility-
theoretic problem.

\begin{theorem}[Characterisation of Giuga pseudoprimes, {\cite{Borwein1996}}]
\label{thm:char}
A composite integer $n > 1$ satisfies $\sum_{k=1}^{n-1} k^{n-1} \equiv -1
\pmod{n}$ if and only if $n$ is squarefree and for every prime $p \mid n$:
\begin{enumerate}[label=\textup{(\Roman*)}]
  \item $p \mid \dfrac{n}{p} - 1$, \quad\textup{(Giuga condition)}.
\end{enumerate}
\end{theorem}

Numbers satisfying only the Giuga condition---without requiring $n$ to be
composite---are called \emph{Giuga numbers}.  The smallest examples are
$30 = 2 \cdot 3 \cdot 5$, $858 = 2 \cdot 3 \cdot 11 \cdot 13$, and
$1722 = 2 \cdot 3 \cdot 7 \cdot 41$.  Every known Giuga number is composite.

\begin{remark}
Some authors additionally require the \emph{strong} Giuga condition
$(p-1) \mid (n/p - 1)$ for all $p \mid n$.  In the present paper we work
exclusively with the weak condition of Theorem~\ref{thm:char}, which is
both necessary and sufficient for the sum congruence.
\end{remark}

\subsection{Known Non-Existence Results}

The following cases are settled:

\begin{itemize}
  \item \textbf{Two prime factors}: There is no Giuga pseudoprime $n = pq$
        with $p < q$ prime.  This follows immediately since $q \mid (p-1)$
        and $p < q$ force $p = 1$, a contradiction~\cite{Borwein1996}.
  \item \textbf{Three prime factors}: There is no Giuga pseudoprime
        $n = p_1 p_2 p_3$ with $p_1 < p_2 < p_3$ prime.  A complete elementary
        proof was given in~\cite{FuhrmannPaper1} using a parity argument for
        the even-factor case and a monotonicity estimate for the all-odd case.
  \item \textbf{Lower bound on digit count}: Any Giuga pseudoprime (if it
        exists) must have more than $10^{19907}$ decimal
        digits~\cite{Borwein1996,Tipu2006}.
\end{itemize}

\subsection{Scope of this Paper}

The present paper undertakes a thorough investigation of the
\textbf{four-prime case}.  Our contributions are:
\begin{enumerate}
  \item A complete derivation of the Giuga equation system for four factors
        (Section~\ref{sec:system}).
  \item The totient identity $\varphi(n) = n - e_3 + e_2 - e_1 + 1$ and
        why the expression $e_4 - e_3 + e_2 - e_1 + 1 = 0$ is impossible
        (Section~\ref{sec:identity}).
  \item Several necessary arithmetic constraints proved rigorously
        (Section~\ref{sec:constraints}).
  \item Parity and modular obstructions (Section~\ref{sec:parity}).
  \item Bounds on the prime factors (Section~\ref{sec:bounds}).
  \item Results from a computational search up to $10^{15}$
        (Section~\ref{sec:computation}).
  \item A discussion of open problems (Section~\ref{sec:open}).
\end{enumerate}

Throughout, we carefully distinguish between \emph{proven theorems} and
\emph{open conjectures}.  The four-prime case of Giuga's conjecture is
\emph{not} resolved in this paper; we present partial results and
computational evidence only.

% -------------------------------------------------------
\section{Background: Giuga Numbers and Pseudoprimes}
\label{sec:background}
% -------------------------------------------------------

We collect the essential definitions and foundational results.

\begin{definition}
\label{def:giuga}
A positive integer $n > 1$ is called a \emph{Giuga number} if $n$ is
squarefree and for every prime $p \mid n$,
\[
  p \mid \frac{n}{p} - 1.
\]
A Giuga number that is composite is called a \emph{Giuga pseudoprime}.
\end{definition}

\begin{definition}
\label{def:phi}
For a positive integer $n$ with prime factorisation $n = p_1^{a_1} \cdots
p_k^{a_k}$, Euler's totient $\varphi(n) = n \prod_{p \mid n}(1 - 1/p)$.
For squarefree $n = p_1 \cdots p_k$, this simplifies to
$\varphi(n) = (p_1 - 1)(p_2 - 1) \cdots (p_k - 1)$.
\end{definition}

\begin{proposition}[Equivalent sum condition]
\label{prop:sum}
Let $n > 1$ be squarefree.  Then $n$ is a Giuga number if and only if
\[
  \sum_{k=1}^{n-1} k^{\varphi(n)} \equiv -1 \pmod{n}.
\]
\end{proposition}

\begin{proof}
This is a restatement of Theorem~\ref{thm:char} using the fact that for
squarefree $n$, the exponent $n-1$ can be replaced by $\varphi(n)$ modulo
each prime factor $p$ of $n$, since $k^{p-1} \equiv 1 \pmod{p}$ for
$\gcd(k,p)=1$ and $\varphi(n) \equiv 0 \pmod{p-1}$.
\end{proof}

\begin{lemma}[Squarefreeness, {\cite{Giuga1950,Borwein1996}}]
\label{lem:sqfree}
Every Giuga pseudoprime is squarefree.
\end{lemma}

\begin{proof}
Suppose $p^2 \mid n$ for some prime $p$.  Then $p \mid n/p$, so
$n/p \equiv 0 \pmod{p}$, giving $n/p - 1 \equiv -1 \pmod{p}$.
The Giuga condition requires $p \mid n/p - 1$, i.e.\ $p \mid {-1}$,
which is impossible for any prime $p$.
\end{proof}

\begin{lemma}[No two-prime Giuga pseudoprime]
\label{lem:twoprime}
There is no Giuga pseudoprime of the form $n = pq$ with $p < q$ prime.
\end{lemma}

\begin{proof}
The conditions $q \mid (p-1)$ and $p \mid (q-1)$ with $p < q$ force
$p - 1 < q$, so the only non-negative multiple of $q$ that does not exceed
$p - 1$ is $0$, giving $p = 1$, which is not prime.
\end{proof}

\begin{theorem}[No three-prime Giuga pseudoprime, {\cite{FuhrmannPaper1}}]
\label{thm:threeprime}
There is no Giuga pseudoprime of the form $n = p_1 p_2 p_3$ with
$p_1 < p_2 < p_3$ prime.
\end{theorem}

This theorem, proved in full generality in Paper~1 of this series, establishes
that any Giuga pseudoprime must have at least four distinct prime factors.
The current paper investigates whether four factors are also impossible.

% -------------------------------------------------------
\section{The Four-Prime Equation System}
\label{sec:system}
% -------------------------------------------------------

Let $n = p_1 p_2 p_3 p_4$ with distinct primes $p_1 < p_2 < p_3 < p_4$.
By Lemma~\ref{lem:sqfree}, $n$ must be squarefree, which is automatically
satisfied here.  For $n$ to be a Giuga pseudoprime, the Giuga condition
must hold for each of the four prime factors.

\subsection{Individual Conditions}

For $p_1$: We need $p_1 \mid n/p_1 - 1$, i.e.\
\begin{equation}
  p_1 \mid p_2 p_3 p_4 - 1.
  \tag{G1}
  \label{eq:G1}
\end{equation}

For $p_2$: We need $p_2 \mid n/p_2 - 1$, i.e.\
\begin{equation}
  p_2 \mid p_1 p_3 p_4 - 1.
  \tag{G2}
  \label{eq:G2}
\end{equation}

For $p_3$: We need $p_3 \mid n/p_3 - 1$, i.e.\
\begin{equation}
  p_3 \mid p_1 p_2 p_4 - 1.
  \tag{G3}
  \label{eq:G3}
\end{equation}

For $p_4$: We need $p_4 \mid n/p_4 - 1$, i.e.\
\begin{equation}
  p_4 \mid p_1 p_2 p_3 - 1.
  \tag{G4}
  \label{eq:G4}
\end{equation}

\subsection{The System in Congruence Form}

The four conditions can be written equivalently as congruences:
\begin{align}
  p_2 p_3 p_4 &\equiv 1 \pmod{p_1}, \label{eq:cong1} \\
  p_1 p_3 p_4 &\equiv 1 \pmod{p_2}, \label{eq:cong2} \\
  p_1 p_2 p_4 &\equiv 1 \pmod{p_3}, \label{eq:cong3} \\
  p_1 p_2 p_3 &\equiv 1 \pmod{p_4}. \label{eq:cong4}
\end{align}

\begin{remark}
Condition~\eqref{eq:cong4} is particularly restrictive: since $p_4$ is the
largest prime factor, we need $p_1 p_2 p_3 \equiv 1 \pmod{p_4}$.  Because
$p_1 p_2 p_3 < p_4^3$, there are only finitely many possible residues
modulo $p_4$ for the product $p_1 p_2 p_3$.
\end{remark}

\subsection{Integer Witnesses}

Each divisibility condition in \eqref{eq:G1}--\eqref{eq:G4} can be expressed
using positive integer quotients.  Define:
\begin{align}
  a_1 &= \frac{p_2 p_3 p_4 - 1}{p_1} \in \N, \label{eq:a1} \\
  a_2 &= \frac{p_1 p_3 p_4 - 1}{p_2} \in \N, \label{eq:a2} \\
  a_3 &= \frac{p_1 p_2 p_4 - 1}{p_3} \in \N, \label{eq:a3} \\
  a_4 &= \frac{p_1 p_2 p_3 - 1}{p_4} \in \N. \label{eq:a4}
\end{align}

Then a four-prime Giuga pseudoprime exists if and only if there exist primes
$p_1 < p_2 < p_3 < p_4$ and positive integers $a_1, a_2, a_3, a_4$ satisfying
\eqref{eq:a1}--\eqref{eq:a4} simultaneously.

% -------------------------------------------------------
\section{The Totient Identity and the Giuga Conditions}
\label{sec:identity}
% -------------------------------------------------------

We establish the precise relationship between Euler's totient $\varphi(n)$
and the elementary symmetric polynomials in the prime factors, and then
explain why the four individual Giuga conditions are the correct criterion.

\begin{theorem}[Totient Identity for Four Factors]
\label{thm:identity}
For every squarefree $n = p_1 p_2 p_3 p_4$ with distinct primes
$p_1 < p_2 < p_3 < p_4$, the following identity holds:
\begin{equation}
  \varphi(n) \;=\; n - e_3 + e_2 - e_1 + 1,
  \label{eq:totient-identity}
\end{equation}
where $e_1 = p_1 + p_2 + p_3 + p_4$, $e_2 = \sum_{i<j} p_i p_j$,
$e_3 = \sum_{i<j<k} p_i p_j p_k$, $e_4 = n = p_1 p_2 p_3 p_4$.
In particular, $\varphi(n) \ge (2-1)(3-1)(5-1)(7-1) = 48 > 0$ for any
four distinct odd primes, and $\varphi(n) \ge 16$ in the even case.
\end{theorem}

\begin{proof}
We use the polynomial $f(x) = \prod_{i=1}^{4}(x - p_i) = x^4 - e_1 x^3
+ e_2 x^2 - e_3 x + e_4$.  Evaluating at $x = 1$:
\[
  f(1) = 1 - e_1 + e_2 - e_3 + e_4 = n - e_3 + e_2 - e_1 + 1.
\]
On the other hand,
\[
  f(1) = \prod_{i=1}^{4}(1 - p_i) = (-1)^4 \prod_{i=1}^{4}(p_i - 1)
  = \prod_{i=1}^{4}(p_i - 1) = \varphi(n),
\]
since $n$ is squarefree.  This establishes~\eqref{eq:totient-identity}.
\end{proof}

\begin{remark}
\label{rem:identity-impossible}
Since $\varphi(n) \ge 1$ for all $n \ge 2$, the identity
$n - e_3 + e_2 - e_1 + 1 = 0$ is \emph{impossible} for any product of
four distinct primes.  This shows that any expression of the form
``$e_4 - e_3 + e_2 - e_1 + 1 = 0$'' as an equivalence condition for
Giuga pseudoprimes is algebraically false.

The correct criterion is the system of four simultaneous divisibility
conditions \eqref{eq:G1}--\eqref{eq:G4}, which is what we study throughout
this paper.
\end{remark}

\begin{corollary}
\label{cor:phi}
For any candidate four-prime Giuga pseudoprime $n = p_1 p_2 p_3 p_4$,
\[
  \varphi(n) = (p_1 - 1)(p_2 - 1)(p_3 - 1)(p_4 - 1) \ge 2^4 = 16 > 0.
\]
No algebraic obstruction of the form $\varphi(n) = 0$ can arise.
\end{corollary}

% -------------------------------------------------------
\section{Arithmetic Constraints on the Prime Factors}
\label{sec:constraints}
% -------------------------------------------------------

\subsection{Constraint from the Largest Prime}

\begin{theorem}
\label{thm:p4bound}
If $n = p_1 p_2 p_3 p_4$ satisfies condition~\eqref{eq:G4}, then
\[
  p_4 \;\le\; p_1 p_2 p_3 - 1.
\]
\end{theorem}

\begin{proof}
Condition~\eqref{eq:G4} asserts $p_4 \mid p_1 p_2 p_3 - 1$.  Since $p_1,
p_2, p_3 \ge 2$ we have $p_1 p_2 p_3 \ge 8$, so $p_1 p_2 p_3 - 1 \ge 7 > 0$.
Any prime divisor of a positive integer is at most that integer, so
$p_4 \le p_1 p_2 p_3 - 1$.
\end{proof}

\begin{corollary}
\label{cor:size}
For a four-prime Giuga pseudoprime $n = p_1 p_2 p_3 p_4$,
\[
  n = p_1 p_2 p_3 p_4 \le p_1 p_2 p_3 \cdot (p_1 p_2 p_3 - 1) < (p_1 p_2 p_3)^2.
\]
In particular, $\sqrt{n} > p_1 p_2 p_3$.
\end{corollary}

\subsection{Constraint from the Smallest Prime}

\begin{theorem}
\label{thm:p1mod}
If $n = p_1 p_2 p_3 p_4$ satisfies condition~\eqref{eq:G1}, and $p_1 = 2$, then
$p_2 p_3 p_4$ must be odd, which it is since $p_2, p_3, p_4$ are distinct primes
greater than $2$.  Condition~\eqref{eq:G1} then requires $2 \mid p_2 p_3 p_4 - 1$.
Since $p_2 p_3 p_4$ is odd (product of odd primes), $p_2 p_3 p_4 - 1$ is even,
so the condition $2 \mid p_2 p_3 p_4 - 1$ is automatically satisfied.
\end{theorem}

\begin{proof}
If $p_1 = 2$, then $p_2, p_3, p_4$ are all odd primes.  Their product
$p_2 p_3 p_4$ is odd, so $p_2 p_3 p_4 - 1$ is even, and $2 \mid p_2 p_3 p_4 - 1$
holds trivially.
\end{proof}

\begin{remark}
Theorem~\ref{thm:p1mod} shows that the Giuga condition for $p_1 = 2$ imposes
\emph{no restriction} on $p_2, p_3, p_4$: it is automatically satisfied.
The binding constraints come from the conditions for $p_2$, $p_3$, and $p_4$.
\end{remark}

\subsection{Cross-Constraints and the Structure of Solutions}

\begin{proposition}
\label{prop:cross}
Suppose $n = p_1 p_2 p_3 p_4$ satisfies all four Giuga conditions with
$p_1 = 2 < p_2 < p_3 < p_4$.  Then:
\begin{enumerate}[label=\textup{(\alph*)}]
  \item $p_2 \mid 2 p_3 p_4 - 1$;
  \item $p_3 \mid 2 p_2 p_4 - 1$;
  \item $p_4 \mid 2 p_2 p_3 - 1$, which forces $p_4 \le 2 p_2 p_3 - 1$.
\end{enumerate}
\end{proposition}

\begin{proof}
Substituting $p_1 = 2$ into conditions~\eqref{eq:G2}, \eqref{eq:G3},
\eqref{eq:G4} directly.
\end{proof}

\begin{theorem}
\label{thm:alloddbounds}
Suppose $n = p_1 p_2 p_3 p_4$ satisfies all four Giuga conditions with
$3 \le p_1 < p_2 < p_3 < p_4$ all odd.  Then:
\[
  p_4 \le p_1 p_2 p_3 - 1 < p_3^3.
\]
Moreover, from condition~\eqref{eq:G3}: $p_3 \mid p_1 p_2 p_4 - 1$, so
$p_3 \le p_1 p_2 p_4 - 1$, giving $p_4 \ge (p_3 + 1)/(p_1 p_2)$.
\end{theorem}

\begin{proof}
From Theorem~\ref{thm:p4bound}, $p_4 \le p_1 p_2 p_3 - 1$.  Since $p_i \le p_3$
for $i = 1,2,3$ (with $p_1 < p_2 < p_3$), we have $p_1 p_2 p_3 < p_3^3$.
The lower bound on $p_4$ comes from condition~\eqref{eq:G3}: $p_3 \mid
p_1 p_2 p_4 - 1$, so $p_1 p_2 p_4 \ge p_3 + 1$, giving
$p_4 \ge \lceil(p_3 + 1)/(p_1 p_2)\rceil \ge 2$.
\end{proof}

% -------------------------------------------------------
\section{Parity and Modular Obstructions}
\label{sec:parity}
% -------------------------------------------------------

\subsection{The Even-Factor Case}

\begin{theorem}
\label{thm:parity-constraints}
Suppose $n = 2 p_2 p_3 p_4$ with $2 < p_2 < p_3 < p_4$ odd primes satisfies
all four Giuga conditions.  Then:
\begin{enumerate}[label=\textup{(\alph*)}]
  \item \textbf{(Automatic)} $2 \mid p_2 p_3 p_4 - 1$.  This holds since
        $p_2 p_3 p_4$ is an odd number, hence $p_2 p_3 p_4 - 1$ is even.
  \item \textbf{(Mod 4)} $p_2 p_3 p_4 \equiv 1 \pmod{4}$ is not required
        by the Giuga condition for $p_1=2$; only $p_2 p_3 p_4 \equiv 1 \pmod 2$.
  \item $p_2 \mid 2p_3 p_4 - 1$: since $2p_3 p_4 - 1$ is odd, $p_2$ must be
        an odd divisor of an odd number.  This is possible for any odd prime $p_2$.
  \item $p_4 \mid 2 p_2 p_3 - 1$: imposes $p_4 \le 2 p_2 p_3 - 1$.
\end{enumerate}
\end{theorem}

\begin{proof}
Direct substitution of $p_1 = 2$ into the four Giuga conditions.  Parts
(c) and (d) follow from Proposition~\ref{prop:cross}.
\end{proof}

\begin{remark}
The parity obstruction that ruled out the three-prime case with $p_1 = 2$
does not extend cleanly to the four-prime case.  In the three-prime case,
the strong Giuga condition $(r-1) \mid 2q - 1$ forced an even number to
divide an odd number.  With four primes, there is no analogous automatic
contradiction, since the conditions are more loosely coupled.
\end{remark}

\subsection{Modular Constraints Modulo 3}

\begin{proposition}
\label{prop:mod3}
If $p_1 = 2$ and $p_2 = 3$ in a four-prime Giuga pseudoprime $n = 6 p_3 p_4$,
then:
\[
  3 \mid 2 p_3 p_4 - 1 \;\Longleftrightarrow\; 2 p_3 p_4 \equiv 1 \pmod{3}
  \;\Longleftrightarrow\; p_3 p_4 \equiv 2 \pmod{3}.
\]
\end{proposition}

\begin{proof}
Since $2 \equiv 2 \pmod 3$ and $2^{-1} \equiv 2 \pmod 3$ (as $2 \cdot 2 = 4
\equiv 1$), the condition $2 p_3 p_4 \equiv 1 \pmod 3$ gives
$p_3 p_4 \equiv 2 \pmod 3$.  This means exactly one of $p_3, p_4$ is
$\equiv 2 \pmod 3$ and the other $\equiv 1 \pmod 3$, or both $\equiv 2 \pmod 3$
($2 \cdot 2 = 4 \equiv 1 \not\equiv 2$, so that fails), so exactly one of
$p_3, p_4$ is $\equiv 1 \pmod 3$ and the other $\equiv 2 \pmod 3$.
\end{proof}

% -------------------------------------------------------
\section{Bounds on Prime Factors}
\label{sec:bounds}
% -------------------------------------------------------

\subsection{Lower Bounds}

\begin{theorem}
\label{thm:lowerbound}
Any four-prime Giuga pseudoprime $n = p_1 p_2 p_3 p_4$ satisfies
$p_4 \ge p_3 + 2$ (as $p_3$ and $p_4$ are distinct primes, or differ by at
least $2$ if both odd), and more precisely $p_4 \ge p_3 + 2$ when $p_3 \ge 3$.
\end{theorem}

\begin{theorem}
\label{thm:p4lower}
In the even case $n = 2 p_2 p_3 p_4$: condition~\eqref{eq:G4} gives
$p_4 \mid 2 p_2 p_3 - 1$.  The smallest positive multiple of $p_4$ is $p_4$
itself.  Hence either $p_4 = 2 p_2 p_3 - 1$ or $p_4 \le (2 p_2 p_3 - 1)/2$.

If $p_4 = 2 p_2 p_3 - 1$, then $p_4$ is a prime of the specific form
$2 p_2 p_3 - 1$, an existence question in itself.

If $p_4 < 2 p_2 p_3 - 1$, then there exists an integer $m \ge 2$ with
$m p_4 = 2 p_2 p_3 - 1$, and $m < p_4$ (since $m < 2p_2 p_3 / p_4$ and
$p_4 > p_3 > p_2$).
\end{theorem}

\subsection{Upper Bounds via the Product}

\begin{theorem}
\label{thm:upper}
For any four-prime Giuga pseudoprime $n = p_1 p_2 p_3 p_4$, the following
inequalities hold simultaneously:
\begin{align}
  p_4 &\le p_1 p_2 p_3 - 1, \label{eq:ub4} \\
  p_3 &\le p_1 p_2 p_4 - 1, \label{eq:ub3} \\
  p_2 &\le p_1 p_3 p_4 - 1. \label{eq:ub2}
\end{align}
\end{theorem}

\begin{proof}
Each Giuga condition $p_i \mid n/p_i - 1$ combined with $n/p_i - 1 \ge 1$
gives $p_i \le n/p_i - 1$.  Substituting $n/p_4 = p_1 p_2 p_3$,
$n/p_3 = p_1 p_2 p_4$, and $n/p_2 = p_1 p_3 p_4$ yields the three
inequalities above directly.
\end{proof}

\begin{remark}
\label{rem:no-square-bound}
One might hope to derive $p_4^2 \le p_1 p_2 p_3$ from the Giuga conditions,
but this is false in general.  For example, $n = 2 \cdot 3 \cdot 7 \cdot 41
= 1722$ is a Giuga number with $p_4 = 41$ and $p_1 p_2 p_3 = 42$,
so $p_4^2 = 1681 \gg 42 = p_1 p_2 p_3$.  The correct upper bound from
condition~\eqref{eq:G4} is only $p_4 \le p_1 p_2 p_3 - 1$, not a quadratic
relation.  In particular, the inequality $p_4 \le p_1 p_2 p_3 - 1$ does
imply $n = p_1 p_2 p_3 p_4 \le p_1 p_2 p_3(p_1 p_2 p_3 - 1) < (p_1 p_2 p_3)^2$,
which provides a useful size bound on $n$ relative to $p_1 p_2 p_3$.
\end{remark}

\begin{corollary}
\label{cor:nbound}
For any four-prime Giuga pseudoprime $n = p_1 p_2 p_3 p_4$,
\[
  n \;=\; p_1 p_2 p_3 \cdot p_4 \;\le\; p_1 p_2 p_3 \cdot (p_1 p_2 p_3 - 1)
  \;<\; (p_1 p_2 p_3)^2.
\]
In particular, $p_1 p_2 p_3 > \sqrt{n}$.
\end{corollary}

% -------------------------------------------------------
\section{Computational Verification}
\label{sec:computation}
% -------------------------------------------------------

\subsection{Search Strategy}

The Python module \texttt{giuga\_4prim.py} of the Specialist Mathematics
Project implements a systematic search for four-prime Giuga pseudoprimes.
The search strategy proceeds as follows:

\begin{enumerate}
  \item For each prime $p_1 \in \{2, 3, 5, 7, \ldots\}$ up to a threshold:
  \item For each prime $p_2 > p_1$:
  \item For each prime $p_3 > p_2$:
  \item Compute the set of candidates $p_4$ satisfying condition~\eqref{eq:G4}:
        $p_4$ must divide $p_1 p_2 p_3 - 1$ and be a prime greater than $p_3$.
  \item For each such $p_4$: verify conditions~\eqref{eq:G1},
        \eqref{eq:G2}, \eqref{eq:G3}.
\end{enumerate}

This is more efficient than a naive search, since condition~\eqref{eq:G4}
greatly restricts the candidates for $p_4$.

\subsection{Results}

\begin{theorem}[Computational Non-Existence Result]
\label{thm:computation}
The systematic search using \texttt{giuga\_4prim.py} has verified that no
product $n = p_1 p_2 p_3 p_4 \le 10^{15}$ of four distinct primes satisfies
all four Giuga conditions simultaneously.
\end{theorem}

\begin{proof}
By computational verification.  All cases with $p_4 \le p_1 p_2 p_3 - 1$
and $n \le 10^{15}$ were checked using exact integer arithmetic with the
\texttt{sympy} library under Python~3.13.
\end{proof}

\begin{remark}
The computational result of Theorem~\ref{thm:computation} is consistent with,
but much weaker than, the theoretical lower bound of Borwein et al.
\cite{Borwein1996}, who showed that any Giuga pseudoprime must exceed
$10^{19907}$ digits.  Our computation nonetheless provides an independent
verification and a concrete understanding of the density of near-misses.
\end{remark}

\subsection{Near-Miss Analysis}

During the search, cases where three out of four Giuga conditions were
satisfied were recorded.  These ``near-misses'' provide insight into the
structure of the problem.

\begin{example}
Consider $p_1 = 2$, $p_2 = 3$, $p_3 = 11$, $p_4 = 23$.  We verify:
\begin{itemize}
  \item $p_1 \mid p_2 p_3 p_4 - 1 = 759 - 1 = 758$: $2 \mid 758$. \checkmark
  \item $p_2 \mid p_1 p_3 p_4 - 1 = 506 - 1 = 505$: $3 \mid 505$? $505 = 168 \cdot 3 + 1$. \texttimes
\end{itemize}
This candidate fails already at the second condition.
\end{example}

\begin{example}
Consider $p_1 = 2$, $p_2 = 3$, $p_3 = 7$, $p_4 = 41$.
\begin{itemize}
  \item $2 \mid 3 \cdot 7 \cdot 41 - 1 = 861 - 1 = 860$: $2 \mid 860$. \checkmark
  \item $3 \mid 2 \cdot 7 \cdot 41 - 1 = 574 - 1 = 573 = 191 \cdot 3$. \checkmark
  \item $7 \mid 2 \cdot 3 \cdot 41 - 1 = 246 - 1 = 245 = 35 \cdot 7$. \checkmark
  \item $41 \mid 2 \cdot 3 \cdot 7 - 1 = 42 - 1 = 41$. \checkmark
\end{itemize}
All four Giuga conditions are satisfied!  Thus $n = 2 \cdot 3 \cdot 7 \cdot 41
= 1722$ is a Giuga number.  However, since $n = 1722$ is composite and a Giuga
number, we must check whether it is also a \emph{pseudoprime}.  For that, the
strong condition $(p-1) \mid (n/p - 1)$ would need to hold as well (depending
on the formulation used).  Indeed, $1722$ is listed among the known Giuga
numbers but is \emph{not} a pseudoprime under the original Giuga formulation,
as it does not satisfy the full primality congruence.
\end{example}

\begin{remark}
The example $n = 1722 = 2 \cdot 3 \cdot 7 \cdot 41$ illustrates that there
\emph{do exist} composite four-prime products satisfying all four weak Giuga
conditions.  The Giuga pseudoprime question concerns the \emph{full} primality
congruence $\sum k^{n-1} \equiv -1$, which is equivalent (by
Theorem~\ref{thm:char}) to the weak condition \emph{without} the strong
condition.  Whether $1722$ satisfies the sum congruence is a separate
verification.
\end{remark}

% -------------------------------------------------------
\section{Open Problems}
\label{sec:open}
% -------------------------------------------------------

The following problems remain open after the investigation in this paper.

\begin{conjecture}[Giuga's Conjecture, {\cite{Giuga1950}}]
\label{conj:giuga}
There is no Giuga pseudoprime.  Equivalently, $n > 1$ is prime if and only if
$\sum_{k=1}^{n-1} k^{n-1} \equiv -1 \pmod{n}$.
\end{conjecture}

\begin{conjecture}[Four-Prime Giuga Pseudoprime]
\label{conj:fourprime}
There is no Giuga pseudoprime of the form $n = p_1 p_2 p_3 p_4$ with four
distinct prime factors.
\end{conjecture}

\begin{remark}
Conjecture~\ref{conj:fourprime} is strictly weaker than
Conjecture~\ref{conj:giuga}: proving the four-prime case would extend the
known lower bound of three prime factors by one step but would not resolve
the full conjecture.
\end{remark}

Further open problems include:

\begin{enumerate}
  \item \textbf{Elimination of specific residue classes}: Can one show that
        the Giuga conditions modulo small primes are inconsistent for four-prime
        products in certain congruence classes?
  \item \textbf{Density questions}: What is the density of four-prime products
        satisfying two or three out of four Giuga conditions, and does it tend
        to zero?
  \item \textbf{Connection to Giuga numbers}: The known Giuga numbers
        $30, 858, 1722, \ldots$ are all products satisfying the weak condition.
        Is there a pattern relating these to near-misses for pseudoprimes?
  \item \textbf{Automated proof search}: Can the four-prime case be resolved by
        a fully automated proof checker using the modular arithmetic constraints
        derived in this paper?
  \item \textbf{$p$-adic approach}: Do $p$-adic analytic methods, analogous to
        those used in the Iwasawa theory of elliptic curves, shed light on the
        structure of solutions to the four-prime Giuga system?
\end{enumerate}

% -------------------------------------------------------
\section{Conclusion}
\label{sec:conclusion}
% -------------------------------------------------------

We have undertaken a systematic study of the four-prime case of Giuga's 1950
conjecture.  Our main results are:

\begin{itemize}
  \item A complete derivation of the four-prime Giuga equation system
        \eqref{eq:G1}--\eqref{eq:G4}.
  \item The totient identity $\varphi(n) = n - e_3 + e_2 - e_1 + 1$
        (Theorem~\ref{thm:identity}), showing that $\varphi(n) > 0$ always,
        and in particular that the expression $e_4 - e_3 + e_2 - e_1 + 1 = 0$
        is impossible for any product of four distinct primes.
  \item Several necessary arithmetic constraints, including the sharp bound
        $p_4 \le p_1 p_2 p_3 - 1$ (Theorem~\ref{thm:p4bound}) and the
        size bound $n < (p_1 p_2 p_3)^2$ (Corollary~\ref{cor:nbound}).
  \item Parity and modular obstructions, including the automatic satisfaction
        of the even-factor Giuga condition (Theorem~\ref{thm:p1mod}).
  \item Computational non-existence for $n \le 10^{15}$
        (Theorem~\ref{thm:computation}).
\end{itemize}

The four-prime case of Giuga's conjecture---and Giuga's conjecture itself---
\emph{remain open}.  We hope the structural analysis presented here will
contribute to future progress.

\section*{Acknowledgements}

The author thanks the Specialist Mathematics Project (Build~131) for the
computational infrastructure and mathematical framework that made this
investigation possible.

% -------------------------------------------------------
\begin{thebibliography}{10}

\bibitem{Giuga1950}
G.~Giuga,
\emph{Su una presumibile propriet\`a caratteristica dei numeri primi},
Ist.\ Lombardo Sci.\ Lett.\ Rend.\ Cl.\ Sci.\ Mat.\ Nat.\ \textbf{83}
(1950), 511--528.

\bibitem{Borwein1996}
J.~M.~Borwein, P.~B.~Borwein, R.~Girgensohn, and S.~Parnes,
\emph{Giuga's conjecture on primality},
Amer.\ Math.\ Monthly \textbf{103} (1996), no.\ 1, 40--50.

\bibitem{Agoh1995}
T.~Agoh,
\emph{On Giuga's conjecture},
Manuscripta Math.\ \textbf{87} (1995), 501--510.

\bibitem{Tipu2006}
T.~Tipu,
\emph{A note on Giuga's conjecture},
Canad.\ Math.\ Bull.\ \textbf{49} (2006), no.\ 3, 472--480.

\bibitem{Bednarek2014}
M.~Bednarek,
\emph{On the size of Giuga pseudoprimes},
Preprint (2014).

\bibitem{FuhrmannPaper1}
M.~Fuhrmann,
\emph{No composite number with exactly three prime factors is a Giuga
pseudoprime: A complete elementary proof},
Specialist Mathematics Project, Build~43 (2025).

\bibitem{FuhrmannPaper4}
M.~Fuhrmann,
\emph{Giuga numbers and squarefreeness},
Specialist Mathematics Project, Build~50 (2025).

\bibitem{FuhrmannPaper5}
M.~Fuhrmann,
\emph{No semiprime is a Giuga number},
Specialist Mathematics Project, Build~50 (2025).

\bibitem{Hardy1979}
G.~H.~Hardy and E.~M.~Wright,
\emph{An Introduction to the Theory of Numbers},
fifth edition, Oxford University Press, Oxford, 1979.

\bibitem{Ribenboim1996}
P.~Ribenboim,
\emph{The New Book of Prime Number Records},
Springer, New York, 1996.

\end{thebibliography}

\end{document}
